\documentclass[11pt,a4paper]{article}

\usepackage[utf8]{inputenc}
\usepackage[T1]{fontenc}
\usepackage{geometry}
\usepackage{hyperref}
\usepackage{enumitem}
\usepackage{xcolor}
\usepackage{titlesec}

\geometry{margin=2.5cm}

\title{
TWG$\cdot$FAS\\
Manifesto \& Working Mode\\
\vspace{0.3em}
\large Technical Working Group for Formal Automotive Systems
}

\author{Working Draft}
\date{\today}

\begin{document}

\maketitle
\begin{center}
\textbf{Origin:} https://twg-fas.org
\end{center}

\tableofcontents

\newpage

\section{Preamble}

The Technical Working Group for Formal Automotive Systems (TWG$\cdot$FAS)
is an independent, non-profit, vendor-neutral initiative dedicated to
the development of open standards for formal behavioral representation
in automotive systems.

Our objective is to establish openly specified, verifiable and
interoperable foundations for future AI-driven mobility systems.

TWG$\cdot$FAS exists to enable collaboration between industry,
academia, regulators, open-source communities and independent
researchers.

Participation is voluntary and open.

\section{Mission}

The mission of TWG$\cdot$FAS is:

\begin{enumerate}
\item To develop open, publicly available standards.
\item To improve interoperability across automotive ecosystems.
\item To support safety, verification and certification activities.
\item To promote transparent engineering processes.
\item To create vendor-neutral technical foundations.
\item To foster scientific and industrial collaboration.
\end{enumerate}

\section{Core Principles}

\subsection{Openness}

All specifications shall be publicly accessible.

Technical discussions, decisions and working drafts should be published
whenever legally possible.

\subsection{Vendor Neutrality}

No company, institution, individual, government body or technology
provider shall control the direction of the standard.

The standard shall remain independent of any specific product or
commercial implementation.

\subsection{Interoperability}

Standards shall be designed to enable independent implementations and
portable tooling.

A conforming implementation must not require the use of proprietary
technology.

\subsection{Transparency}

Decision processes, voting outcomes and technical rationales shall be
documented.

Technical disagreements are considered valuable engineering input.

\subsection{Consensus}

TWG$\cdot$FAS seeks consensus rather than majority rule wherever
practical.

When consensus cannot be achieved, decisions may be made through a
transparent voting process.

\subsection{Technical Merit}

Arguments shall be evaluated on technical merit, evidence and reasoning
rather than organizational status, seniority or commercial influence.

\section{Ethical Charter}

Members agree to uphold the following principles:

\begin{enumerate}
\item Respectful communication.
\item Evidence-based technical discussion.
\item Honest disclosure of limitations and assumptions.
\item Attribution of ideas and contributions.
\item Responsible handling of safety-critical topics.
\item Rejection of discrimination and harassment.
\item Commitment to constructive collaboration.
\end{enumerate}

The working group shall maintain a professional and welcoming
environment for all participants.

\section{Code of Conduct}

All participants shall:

\begin{itemize}
\item Act professionally.
\item Treat others with respect.
\item Critique ideas rather than people.
\item Avoid personal attacks.
\item Avoid disruptive behavior.
\item Disclose known conflicts of interest.
\item Support an inclusive technical culture.
\end{itemize}

Violations may result in moderation actions or suspension from working
group activities.

\section{Non-Profit Character}

TWG$\cdot$FAS is not established to generate profit.

The initiative exists for:

\begin{itemize}
\item Standardization
\item Research collaboration
\item Knowledge sharing
\item Public benefit
\item Technical advancement
\end{itemize}

Membership shall not provide exclusive commercial advantage.

No participant shall receive privileged access to specifications.

\section{Governance}

\subsection{Founding Members}

Founding members initiate and organize the working group.

Their role is administrative rather than authoritative.

\subsection{Contributors}

Any individual or organization may contribute proposals, comments,
drafts, implementations or review feedback.

\subsection{Editors}

Editors maintain technical documents and working drafts.

Editors do not possess veto authority.

\subsection{Working Groups}

Specialized working groups may be created for topics including:

\begin{itemize}
\item VML Language Design
\item Formal Semantics
\item Verification
\item Runtime Monitoring
\item Conformance Testing
\item Reference Implementations
\end{itemize}

\section{Standards Development Process}

The development lifecycle consists of:

\begin{enumerate}
\item Problem Statement
\item Proposal Draft
\item Public Review
\item Working Draft
\item Candidate Standard
\item Reference Implementations
\item Conformance Validation
\item Standard Publication
\end{enumerate}

All stages should be publicly visible whenever possible.

\section{Intellectual Property Policy}

Contributors retain ownership of their original work.

By contributing to TWG$\cdot$FAS specifications, contributors grant
the right to publish, distribute and maintain contributions as part
of the standardization process.

Contributors are encouraged to disclose patents that may be essential
for implementation.

\section{Licensing Model}

\subsection{Standards and Specifications}

Specification texts shall be published under:

\begin{center}
\textbf{Creative Commons Attribution 4.0 (CC-BY-4.0)}
\end{center}

This permits:

\begin{itemize}
\item commercial use
\item academic use
\item redistribution
\item translation
\item derivative works
\end{itemize}

provided attribution is maintained.

\subsection{Reference Implementations}

Reference software shall preferably use:

\begin{center}
\textbf{Apache License 2.0}
\end{center}

to enable broad industrial and academic adoption.

\subsection{Patent Commitment}

Participants are encouraged to provide royalty-free licenses for
essential claims required to implement the standard.

The goal is unrestricted implementation by any party.

\section{Conformance and Certification}

TWG$\cdot$FAS may publish:

\begin{itemize}
\item Conformance test suites
\item Validation procedures
\item Reference datasets
\item Compliance guidelines
\end{itemize}

Conformance shall be based on objective technical criteria.

\section{Relationship to Regulation}

TWG$\cdot$FAS does not act as a regulatory body.

The group seeks constructive collaboration with regulators,
certification authorities and standardization organizations.

\section{Long-Term Vision}

TWG$\cdot$FAS aims to create a globally recognized open ecosystem for
formal automotive behavior representation.

The initiative believes that future automotive AI systems require
formal, transparent and verifiable interfaces between human intent,
machine reasoning and vehicle behavior.

Open standards are the foundation for interoperability, trust,
scientific progress and long-term technological sustainability.

\section{Manifesto Statement}

\emph{
Natural language belongs to engineering.\\
Formal languages belong to runtime systems.\\
Open standards belong to everyone.
}

\vspace{1em}

TWG$\cdot$FAS is committed to building an open, transparent and
vendor-neutral foundation for the next generation of safe and
verifiable automotive intelligence.

\end{document}